
\subsubsection{2.1 Data Curation Pipelines}

The Data Management for Training Large Language Models survey \citep{zha2023data} provides an overview of quality filtering, deduplication, and toxicity filtering techniques. RedPajama (Together Computer, 2023) and Dolma \citep{soldaini2024dolma} represent state-of-the-art open curation pipelines processing 100B+ tokens. These pipelines report benchmark improvements but vary multiple factors simultaneously.

C4 \citep{raffel2020t5} established perplexity filtering as standard practice but did not disclose specific thresholds or conduct systematic ablation.

\subsubsection{2.2 Perplexity-Based Quality Filtering}

''How to Train Data-Efficient LLMs'' \citep{tirumala2024dataefficient} demonstrates that perplexity filtering improves training efficiency. ''Perplexed by Perplexity'' \citep{abbas2024perplexed} questions whether optimal thresholds vary by domain and model scale. The standard approach uses KenLM 5-gram models trained on Wikipedia following CCNet methodology \citep{wenzek2020ccnet}.

\subsubsection{2.3 Systematic Curation Studies}

The DataComp benchmark \citep{gadre2023datacomp} provides methodological precedent by systematically varying image curation strategies under controlled conditions, demonstrating that intermediate CLIP Score thresholds outperform both extremes.

This work extends the DataComp methodology to LLM pretraining, treating curation parameters as continuous variables and mapping effect surfaces with fixed-token experimental design.
